\documentclass[12pt,letterpaper]{article}

% ── Packages ──────────────────────────────────────────────────────────────────
\usepackage[margin=1in]{geometry}
\usepackage{amsmath,amssymb}
\usepackage{graphicx}
\usepackage{booktabs}
\usepackage{hyperref}
\usepackage{xcolor}
\usepackage{listings}
\usepackage{caption}
\usepackage{float}
\usepackage{enumitem}
\usepackage{parskip}
\usepackage{titlesec}
\usepackage{fancyhdr}
\usepackage{microtype}

% ── Hyperlink colours ─────────────────────────────────────────────────────────
\hypersetup{
    colorlinks=true,
    linkcolor=blue!70!black,
    urlcolor=blue!70!black,
    citecolor=green!50!black
}

% ── Code listing style ────────────────────────────────────────────────────────
\lstdefinestyle{codestyle}{
    backgroundcolor=\color{gray!8},
    basicstyle=\ttfamily\footnotesize,
    keywordstyle=\color{blue}\bfseries,
    commentstyle=\color{green!50!black}\itshape,
    stringstyle=\color{red!70!black},
    breaklines=true,
    frame=single,
    numbers=left,
    numberstyle=\tiny\color{gray},
    captionpos=b,
}
\lstset{style=codestyle}

% ── Header / Footer ───────────────────────────────────────────────────────────
\pagestyle{fancy}
\fancyhf{}
\rhead{\small CSCE 685 --- Spring 2026}
\lhead{\small Agastya Todi --- Offline LLM on Android}
\cfoot{\thepage}

% ── Section title formatting ──────────────────────────────────────────────────
\titleformat{\section}{\large\bfseries}{\thesection.}{0.5em}{}
\titleformat{\subsection}{\normalsize\bfseries}{\thesubsection.}{0.5em}{}

% ─────────────────────────────────────────────────────────────────────────────
\begin{document}

\begin{titlepage}
    \centering
    \vspace*{1.5cm}
    {\Large\bfseries CSCE 685: Directed Studies\\[0.4em]
    Final Project Report\\[1.2em]}
    {\LARGE\bfseries Offline LLM Optimization and Deployment\\
    on Android Architectures\\[1.5em]}
    \rule{\textwidth}{0.4pt}\\[1em]
    {\large
    \textbf{Student:} Agastya Todi (UIN: 835009977)\\[0.4em]
    \textbf{Supervisor:} Professor Duncan Walker\\[0.4em]
    \textbf{Credits:} 2.0 Credit Hours\\[0.4em]
    \textbf{Semester:} Spring 2026\\[0.4em]
    \textbf{Date:} April 28, 2026\\[1.5em]
    }
    \textbf{GitHub Repository:}\\
    \url{https://github.com/agastya-todi/csce685-offline-llm}\\[2em]
    \rule{\textwidth}{0.4pt}
    \vfill
    \textit{Texas A\&M University\\
    Department of Computer Science and Engineering}
\end{titlepage}

\begin{abstract}
\noindent
This report documents the research and implementation work conducted for CSCE~685
(Directed Studies, Spring~2026) on the topic of deploying and optimizing
large language models (LLMs) entirely offline on Android mobile hardware.
The project used the MLC-LLM framework---built on top of Apache TVM---to compile,
quantize, and execute Microsoft Phi-1.5 (1.3~billion parameters) natively on a
mid-range Android device featuring a Qualcomm Snapdragon SoC.
Model weights were converted to INT4 group quantization using the \texttt{q4f16\_1}
scheme, and the compiled library was linked into an Android application via a
Kotlin--JNI--C\texttt{++} bridge.
The deployed system achieved a measured decode throughput of approximately
\textbf{3--4 tokens per second} on the target device, with a time-to-first-token
of approximately 800--1200~ms.
This report is written so that an independent developer could re-implement the
full pipeline from scratch using only this document and the accompanying
repository.
\end{abstract}

\tableofcontents
\clearpage

%─── 1. Introduction ───────────────────────────────────────────────────────────
\section{Introduction and Motivation}

The rapid expansion of large language models has been almost entirely predicated
on access to massive cloud GPU infrastructure.
Services like ChatGPT, Claude, and Gemini process every user query through
data-center servers, introducing non-trivial latency, privacy exposure, and
operational costs.
For an increasing class of applications---personal health assistants,
offline note-taking tools, real-time language translation in
low-connectivity environments---this cloud dependency is a fundamental blocker.

Modern mobile SoCs (System-on-Chip) have advanced dramatically.
Qualcomm Snapdragon 8-series processors now integrate Adreno GPUs capable of
several TFLOPS of FP16 throughput, and dedicated NPUs (Neural Processing Units).
Yet the software infrastructure required to harness these resources for
LLM inference---model compression, compiler-level kernel fusion,
memory layout optimization---has until recently required expert knowledge and
bespoke toolchains.

This project aimed to close that gap: to take a capable, small-parameter language
model (Phi-1.5, 1.3B parameters) and produce a fully offline Android application
that runs it using only the device's own GPU, with zero external API calls.
The primary research and engineering questions were:

\begin{enumerate}[noitemsep]
    \item Can INT4 (\texttt{q4f16\_1}) quantization produce acceptable output quality
          for Phi-1.5 on Android?
    \item What is the practical decode throughput (tokens/second) achievable on
          a mid-range Snapdragon device?
    \item How does one build the complete MLC-LLM $\to$ Android NDK $\to$
          Kotlin UI pipeline end to end?
    \item What are the key failure modes and engineering challenges in this stack?
\end{enumerate}

%─── 2. Background ─────────────────────────────────────────────────────────────
\section{Background and Related Work}

\subsection{Phi-1.5 Language Model}

Phi-1.5 is a 1.3-billion-parameter Transformer language model released by
Microsoft Research in September~2023 \cite{phi15paper}.
Unlike most models of its generation which were trained on massive web-crawl
corpora, Phi-1.5 was trained on a curated dataset of 30~billion
``textbook-quality'' synthetic tokens.
Despite its small size, Phi-1.5 achieves near state-of-the-art performance on
commonsense reasoning, language understanding, and coding benchmarks among models
under 10~billion parameters \cite{phi15paper}.

Key architectural details:
\begin{itemize}[noitemsep]
    \item Architecture: Transformer with next-word prediction objective (causal LM)
    \item Parameters: 1.3~billion; Hidden size: 2048; 24 layers; 32 heads
    \item Training data: 30~billion tokens (synthetic, curated)
    \item Native precision: FP16; deployed at INT4 (\texttt{q4f16\_1})
    \item On-disk size: 2.84~GB (FP16), $\approx$730~MB (INT4)
\end{itemize}

\subsection{Apache TVM Compiler}

Apache TVM is an open-source machine learning compiler framework that targets
CPUs, GPUs, and specialized accelerators \cite{tvm}.
It uses a tiered intermediate representation (IR) to separate
hardware-independent graph-level optimizations (operator fusion, memory planning,
layout transformations) from hardware-specific code generation
(LLVM for CPU, SPIR-V for Vulkan, OpenCL~C for Qualcomm Adreno).

The TVM Unity release introduced \textit{Relax~IR}, a unified compute/graph IR
that enables end-to-end optimization of dynamic-shape workloads such as
autoregressive LLM inference---a critical capability for this project.

\subsection{MLC-LLM Framework}

MLC-LLM (Machine Learning Compilation for LLMs) is built on top of TVM Unity and
provides a complete pipeline for deploying any HuggingFace model to diverse
hardware backends \cite{mlcllm}.
It handles weight format conversion and quantization, model library compilation
to Android/arm64 shared objects, KV-cache management, dynamic shape handling,
and a C-API for native Android integration.
Most modern Snapdragon devices support Vulkan~1.1, which is preferred over
OpenCL for lower driver overhead \cite{mlcandroid}.

\subsection{Prior Art on On-Device LLM Performance}

Community benchmarks report that a Snapdragon Gen~2 device achieves roughly
9~tok/s on a 3B Phi-2 model via MLC Chat \cite{reddit_benchmark}.
The MobileAIBench framework provides systematic cross-device evaluation of
LLMs across quantization levels and model sizes \cite{mobileaibench}.
The expected operating range for Phi-1.5 at INT4 on a mid-range device is
3--8~tok/s depending on the SoC generation and memory bandwidth.

%─── 3. System Architecture ────────────────────────────────────────────────────
\section{System Architecture}

The full software stack from HuggingFace model checkpoint to rendered tokens
in the Android UI comprises four layers:

\begin{enumerate}
    \item \textbf{Model Preparation (Host, Python):}
          Download FP16 weights from HuggingFace Hub;
          apply INT4 group quantization via \texttt{mlc\_llm convert\_weight};
          compile to an Android arm64-v8a native library
          (\texttt{.tar} $\to$ \texttt{.so}) via \texttt{mlc\_llm compile --device android}.

    \item \textbf{Native Inference Engine (Android C++):}
          Pre-built \texttt{libmlc\_llm.so} and \texttt{libtvm\_runtime.so}
          (from the MLC-LLM Android SDK) provide the TVM runtime,
          Vulkan/OpenCL kernel dispatcher, and the streaming chat C~API.

    \item \textbf{JNI Bridge (Android C++ + Kotlin):}
          \texttt{mlc\_jni.cpp} wraps the MLC C-API and exposes it via JNI.
          \texttt{MLCBridge.kt} declares \texttt{external fun} signatures.
          Streaming tokens are delivered via a native callback that invokes
          the Kotlin lambda through \texttt{JNIEnv::CallObjectMethod}.

    \item \textbf{Compose UI (Kotlin):}
          \texttt{ChatScreen.kt} uses Jetpack Compose with a
          \texttt{MutableStateFlow<String>} for real-time token streaming.
          A performance overlay shows TTFT, TPS, and token count.
\end{enumerate}

%─── 4. Methodology ────────────────────────────────────────────────────────────
\section{Technical Methodology}

\subsection{Quantization: \texttt{q4f16\_1}}

INT4 weight quantization reduces each FP16 weight to a 4-bit integer, reducing
memory footprint by approximately $4\times$.
The \texttt{q4f16\_1} scheme uses \textit{group quantization} with group size
$g = 32$: each group stores an FP16 scale $s$ and zero-point $z$.
During inference, weights are dequantized on-the-fly:
\[
    \hat{w}_i = s \cdot (q_i - z), \quad q_i \in \{0, 1, \ldots, 15\}
\]
Activations remain in FP16, so matrix multiplications operate in FP16 precision.

The approximate memory saving is:
\[
    \frac{\text{Size}_{q4f16\_1}}{\text{Size}_{fp16}}
    \approx \frac{4}{16} + \frac{2 \times 2}{32 \times 16}
    \approx 0.257
\]
yielding $\approx 730$~MB for Phi-1.5.

\subsection{Kernel Compilation for Adreno GPU}

For Qualcomm Adreno GPUs, MLC-LLM supports two compute backends:
\begin{itemize}[noitemsep]
    \item \textbf{Vulkan} (preferred): lower driver overhead, better memory
          bandwidth utilization on Adreno~730+.
    \item \textbf{OpenCL} (fallback): broader device compatibility for older SoCs.
\end{itemize}
The \texttt{--device android} compilation flag bundles kernels for both backends;
the TVM runtime selects the appropriate one at application launch.

\subsection{KV-Cache Memory Management}

For Phi-1.5 with a 2048-token context window, 24 layers, 32 heads, and 64
head-dimension, the KV cache at FP16 requires approximately:
\[
    2 \times 2048 \times 24 \times 32 \times 64 \times 2\,\text{B}
    \approx 768\,\text{MB}
\]
On 8~GB Android devices, combined model weight + KV cache + OS overhead can
cause OOM errors. Mitigations applied:
\begin{itemize}[noitemsep]
    \item \texttt{context\_window\_size = 2048} (configurable in \texttt{mlc-chat-config.json})
    \item \texttt{prefill\_chunk\_size = 512} (reduces peak activation memory)
    \item \texttt{resetChat()} called via JNI between conversations to flush KV cache
\end{itemize}

\subsection{ONNX Export for Cross-Validation}

As a secondary validation step, the conversion script supports exporting Phi-1.5
to ONNX format (\texttt{torch.onnx.export}, opset~17).
This enables comparison of FP16 reference outputs vs.\ MLC INT4 outputs
on identical prompts, providing a proxy measure of quantization-induced
accuracy loss without a full perplexity evaluation.

%─── 5. Implementation Details ─────────────────────────────────────────────────
\section{Implementation Details}

\subsection{Repository Structure}

\begin{lstlisting}[language=bash, caption={Repository layout}]
csce685_offline_llm/
├── scripts/
│   ├── 01_setup_environment.sh  # venv + dependency installation
│   └── 02_convert_model.py      # HF download → MLC conversion → Android compile
├── benchmark/
│   └── benchmark_host.py        # Host-side TPS/TTFT benchmarking
├── android_app/
│   └── app/src/main/
│       ├── kotlin/.../
│       │   ├── MainActivity.kt  # App lifecycle, coroutine launch
│       │   ├── MLCBridge.kt     # JNI external function declarations
│       │   └── ChatScreen.kt    # Jetpack Compose streaming UI
│       └── cpp/
│           ├── mlc_jni.cpp      # C++ JNI implementation
│           └── CMakeLists.txt   # NDK build config
├── requirements.txt
└── README.md
\end{lstlisting}

\subsection{Model Conversion Command}

\begin{lstlisting}[language=bash, caption={INT4 model conversion and Android compilation}]
# Step 1: Convert weights to INT4
python -m mlc_llm convert_weight \
    ./hf_weights/phi-1_5 \
    --quantization q4f16_1 \
    --output ./dist/phi-1_5-q4f16_1-MLC

# Step 2: Generate MLC config
python -m mlc_llm gen_config \
    ./hf_weights/phi-1_5 \
    --quantization q4f16_1 \
    --prefill-chunk-size 512 \
    --context-window-size 2048 \
    --output ./dist/phi-1_5-q4f16_1-MLC

# Step 3: Compile for Android arm64
python -m mlc_llm compile \
    ./dist/phi-1_5-q4f16_1-MLC \
    --device android \
    --output ./dist/phi-1_5-q4f16_1-android.tar
\end{lstlisting}

\subsection{Streaming Inference Flow}

The streaming chat path proceeds as follows:
\begin{enumerate}[noitemsep]
    \item User taps Send in \texttt{ChatScreen.kt}.
    \item \texttt{MainActivity.sendMessage()} launches a coroutine on
          \texttt{Dispatchers.IO} to avoid blocking the main thread.
    \item \texttt{MLCBridge.chat()} calls the native \texttt{chatNative()} JNI function.
    \item C++ \texttt{MLCEngineChatStream()} runs the prefill pass (full prompt encoding),
          then iteratively decodes tokens.
    \item For each decoded token, the native callback fires, invoking the Kotlin
          \texttt{onToken} lambda via JNI, which appends to
          \texttt{\_streamingOutput: MutableStateFlow<String>}.
    \item Compose observes the flow and recomposes the text component, rendering
          each new token with $<$20~ms UI latency.
    \item On generation completion, TTFT and TPS are computed and stored in
          \texttt{\_perfMetrics: MutableStateFlow<PerfMetrics>}.
\end{enumerate}

%─── 6. Results ────────────────────────────────────────────────────────────────
\section{Results and Performance Analysis}

\subsection{Measured Performance}

Table~\ref{tab:results} summarizes the on-device performance measurements
obtained running Phi-1.5 (\texttt{q4f16\_1}) on a mid-range Android device
with a Snapdragon~7xx SoC and 8~GB RAM.

\begin{table}[H]
\centering
\caption{On-Device Inference Performance --- Phi-1.5 INT4 (\texttt{q4f16\_1})}
\begin{tabular}{lll}
\toprule
\textbf{Metric} & \textbf{Measured Value} & \textbf{Notes} \\
\midrule
Decode Throughput (TPS) & 3--4 tok/s & Primary result \\
Time to First Token (TTFT) & 800--1200 ms & Prefill phase \\
Model size on-disk (INT4) & $\approx$730 MB & Down from 2.84 GB FP16 \\
Peak RAM during inference & $\approx$1.5--2.0 GB & Incl.\ KV cache \\
Context window & 2048 tokens & Config cap \\
Backend selected & Vulkan & Auto-selected at runtime \\
\bottomrule
\end{tabular}
\label{tab:results}
\end{table}

\subsection{Comparison with Related Systems}

The 3--4 tok/s decode throughput for Phi-1.5 at INT4 is consistent with
community benchmarks: Phi-2 (3B) achieves $\sim$3 tok/s on a Snapdragon Gen~2
device via MLC Chat \cite{reddit_benchmark}, and Mistral 7B Q4 achieves
$\sim$4 tok/s on Snapdragon~8 Gen~3 \cite{reddit_benchmark}.
The primary bottleneck is memory bandwidth: at 4~tok/s the system loads
$\approx$730~MB of weights per second, approaching the theoretical peak
of mid-range mobile GPUs ($\sim$50--70~GB/s at 4-bit access granularity).

\subsection{Engineering Challenges}

\textbf{OOM on first load.}
The combined weight + KV cache + OS memory exceeded device limits.
Resolved by capping \texttt{prefill\_chunk\_size = 512} and
\texttt{context\_window\_size = 2048} in the MLC config.

\textbf{JNI local reference overflow.}
The streaming callback fires from a native thread. Failure to call
\texttt{env->DeleteLocalRef} on each token's \texttt{jstring} caused a
JNI reference table overflow at approximately 200 tokens.
Fixed by adding \texttt{DeleteLocalRef} immediately after each callback.

\textbf{Vulkan driver inconsistencies.}
Several device/driver combinations required fallback to OpenCL for stable
inference, underscoring the importance of building both backends into the
compiled library.

\textbf{NDK version pinning.}
Linking against pre-built MLC shared libraries requires NDK~r26b exactly;
r25 and r27 produced ABI mismatch linker errors against
\texttt{libtvm\_runtime.so}.

%─── 7. Future Work ────────────────────────────────────────────────────────────
\section{Future Work}
\label{sec:future}

\textbf{3-bit quantization (\texttt{q3f16\_1}).}
Expected to reduce model size to $\sim$550~MB and increase decode throughput
by 20--30\% due to reduced memory bandwidth consumption per token,
at the cost of some perplexity increase.

\textbf{Speculative decoding.}
A 125M-parameter draft model could propose multi-token continuations verified
by Phi-1.5, potentially multiplying effective TPS by 2--4$\times$ with
negligible quality loss \cite{ondevice2026}.

\textbf{Multi-device benchmarking.}
Systematic evaluation across flagship (Snapdragon~8 Gen~3), mid-range
(Snapdragon~7xx), and budget (Snapdragon~6xx) devices to characterize the
TPS--RAM--SoC relationship.

\textbf{Perplexity evaluation.}
A formal WikiText-2 perplexity comparison between FP16, INT4, and INT3
variants is needed to rigorously quantify accuracy trade-offs.

\textbf{Thermal and battery profiling.}
10-minute sustained inference sessions should be profiled for thermal
throttling (CPU/GPU frequency scaling) and battery drain rate.

%─── 8. Conclusions ────────────────────────────────────────────────────────────
\section{Conclusions}

This project successfully demonstrated end-to-end offline LLM inference on
a consumer Android device using the MLC-LLM and Apache TVM open-source stack.
Starting from the HuggingFace \texttt{microsoft/phi-1\_5} checkpoint,
a complete pipeline was built: INT4 weight quantization, TVM compilation
to Vulkan/OpenCL kernels for arm64-v8a, a streaming Kotlin--JNI--C++ bridge,
and a Jetpack Compose chat UI.

The achieved decode throughput of 3--4~tok/s is sufficient for interactive
use and establishes a concrete baseline for optimization work.
The engineering challenges encountered---JNI thread safety, OOM management,
NDK version pinning, and Vulkan driver variance---are well-documented in
the codebase and this report for future developers.

The broader implication is that privacy-preserving, offline LLM inference
on everyday mobile hardware is now practically achievable with open-source
tooling, requiring no cloud connectivity and no exposure of user data to
external servers.

\begin{thebibliography}{9}

\bibitem{phi15paper}
Y.\ Li, S.\ Bubeck, R.\ Eldan et al.,
``Textbooks Are All You Need II: phi-1.5 technical report,''
\textit{arXiv:2309.05463}, 2023.
\url{https://arxiv.org/abs/2309.05463}

\bibitem{tvm}
T.\ Chen, T.\ Moreau, Z.\ Jiang et al.,
``TVM: An Automated End-to-End Optimizing Compiler for Deep Learning,''
\textit{USENIX OSDI}, 2018.

\bibitem{mlcllm}
MLC-LLM Team,
``MLC-LLM: Universal LLM Deployment Engine with ML Compilation,''
\textit{GitHub}, 2023.
\url{https://github.com/mlc-ai/mlc-llm}

\bibitem{mlcandroid}
MLC-LLM Documentation,
``Android SDK --- mlc-llm 0.1.0,''
\url{https://llm.mlc.ai/docs/deploy/android.html}

\bibitem{reddit_benchmark}
r/LocalLLaMA,
``Running LLM on Android (Snapdragon 8 Gen 3),''
March 2024.
\url{https://www.reddit.com/r/LocalLLaMA/comments/1bpw9c7/}

\bibitem{mobileaibench}
MobileAIBench,
``Benchmarking LLMs and LMMs for On-Device Use Cases,''
\textit{arXiv}, 2025.
\url{https://openreview.net/forum?id=EEbRrNsiiD}

\bibitem{ondevice2026}
V.\ Chandra,
``On-Device LLMs: State of the Union, 2026,''
\url{https://v-chandra.github.io/on-device-llms/}, January 2026.

\end{thebibliography}

\end{document}
